% ---
% id: transmitter/decoy/req:decoy-emission
% kind: req
% title: "Emission of the tagged ensemble as declared"
% summary: >-
%   What the engine emits from a tagged declaration is the decoy declaration of
%   [[cor:emitted-decoy]] and nothing else, with every declared number
%   unchanged. Its ledger carries A-WCP in the place of A-ENS.
% prominence: standard
% status: implemented
% uses: [cor:emitted-decoy, def:tagged-source]
% mentions: [ass:wcp]
% assumes: []
% refs: []
% provenance: newly-derived
% implements: certiqs_key.transmitter.decoy.emit:emit_decoy_declaration
% lean:
% ---

The source identity holds at the emitted descriptors because the emitted
ensemble is the tagged one with the declared weights
(Corollary~\ref{cor:emitted-decoy}\,(ii)).

\begin{requirement}\label{req:decoy-emission}
When the engine emits an artefact from a \tref{tagged source declaration} over
a \tref{weak-coherent-pulse model}, the artefact is that declaration's decoy
declaration (Corollary~\ref{cor:emitted-decoy}) and nothing else: the engine
declares no retained system, no measurement, no marginal and no certificate,
and it rounds or otherwise alters no declared number. The matched
\tref{sector structure} among the declared data of that corollary is emitted as
the \tref{retained-register sector map} of
Definition~\ref{def:tagged-source}\,(T5), with
the block kind \iface{proof_sector}, the tag that records a decomposition
licensed by a proof, here source replacement, rather than by the device's own
mode structure;
it names Alice's sectors only as copies of the declared sectors $S_t$, so it
declares no retained system. An artefact emitted from a tagged declaration
carries A-WCP and not A-ENS in its assumption ledger
(Assumption~\ref{ass:wcp}).
\end{requirement}
